\documentclass[12pt,a4paper]{article}
\usepackage[utf8]{inputenc}
\usepackage[T1]{fontenc}
\usepackage{geometry}
\usepackage{amsmath, amssymb}
\usepackage{booktabs}
\usepackage{listings}
\usepackage{xcolor}
\usepackage{hyperref}
\usepackage{parskip}

\geometry{margin=1in}

% ── Python Listing Style ──────────────────────────────────────────────────────
\definecolor{codebg}{RGB}{248,248,248}
\definecolor{codegreen}{rgb}{0,0.6,0}
\definecolor{codegray}{rgb}{0.5,0.5,0.5}
\definecolor{codeblue}{rgb}{0,0,0.8}
\definecolor{codeorange}{rgb}{0.8,0.4,0}

\lstdefinestyle{pythonstyle}{
    language=Python,
    backgroundcolor=\color{codebg},
    basicstyle=\ttfamily\footnotesize,
    keywordstyle=\color{codeblue}\bfseries,
    commentstyle=\color{codegreen}\itshape,
    stringstyle=\color{codeorange},
    numbers=left,
    numberstyle=\tiny\color{codegray},
    stepnumber=1,
    numbersep=8pt,
    tabsize=4,
    showstringspaces=false,
    breaklines=true,
    frame=single,
    captionpos=b,
    morekeywords={self, True, False, None, nn, F, torch, np, pd}
}
\lstset{style=pythonstyle}
% ─────────────────────────────────────────────────────────────────────────────

\title{\textbf{Dokumentasi Kode} \\ \large Strategi Portofolio Berbasis Graph Neural Networks (GNN) \\
\small \texttt{strategy\_comparison\_gnn\_only.ipynb}}
\author{Catatan Tesis}
\date{\today}

\begin{document}

\maketitle
\tableofcontents
\newpage

%=============================================================================
\section{Pendahuluan}
%=============================================================================
Dokumen ini fokus pada implementasi dan evaluasi strategi alokasi portofolio menggunakan 
\textbf{GNN Graph Diversification}. Berbeda dengan metode klasik, pendekatan ini 
menggunakan \textit{Graph Convolutional Networks} (GCN) untuk mempelajari dependensi 
antar aset kripto secara non-linier dan membangun graf korelasi prediktif untuk 
pemilihan aset melalui \textit{Maximum Independent Set} (MIS).

Varian yang diuji mencakup:
\begin{enumerate}
    \item \textbf{GNN Graph Diversification (30 Epoch)}: Pelatihan model durasi pendek.
    \item \textbf{GNN Graph Diversification (50 Epoch)}: Pelatihan model durasi menengah.
    \item \textbf{GNN Graph Diversification (100 Epoch)}: Eksperimen konvergensi lebih dalam.
\end{enumerate}

%=============================================================================
\section{Sel 1 -- Import Libraries}
%=============================================================================

\subsection{Deskripsi}
Mengimpor pustaka utama untuk Deep Learning graf dan ekosistem sains data Python.
\begin{lstlisting}[caption={Import libraries}]
import numpy as np
import pandas as pd
import matplotlib.pyplot as plt
import seaborn as sns
import networkx as nx
from scipy.optimize import minimize
from scipy.linalg import eigh
from scipy import stats
import torch
import torch.nn as nn
import torch.nn.functional as F
from torch_geometric.nn import GCNConv
from torch_geometric.data import Data
import warnings
warnings.filterwarnings('ignore')
plt.style.use('seaborn-v0_8-darkgrid')
\end{lstlisting}

%=============================================================================
\section{Sel 2 -- Load \& Prepare Real Crypto Data}
%=============================================================================

\subsection{Deskripsi}
Memuat dataset return harian aset kripto dari file Excel. Pembersihan dilakukan untuk memastikan 
hanya data sinkron yang digunakan. Bagian ini juga mendefinisikan label sub-periode pasar 
(\textit{Bull, Bear, Recovery}) untuk analisis performa yang lebih mendalam.

\subsection{Kode}
\begin{lstlisting}[caption={Load dan preparasi data lengkap}]
EXCEL_FILE = 'crypto_data_real.xlsx'
df_raw = pd.read_excel(EXCEL_FILE, sheet_name='Returns', index_col=0)
df_raw.index = pd.to_datetime(df_raw.index)
df_raw.sort_index(inplace=True)
df_raw = df_raw.drop(columns=['USDT'], errors='ignore')  # stablecoin

# Hanya simpan baris di mana semua aset sudah aktif (return != 0)
mask = (df_raw != 0).all(axis=1)
df_returns = df_raw[mask].copy()

# Sub-period labels
BULL_END  = '2017-12-31'   # Bull: Nov 2017 - Des 2017
BEAR_END  = '2018-12-31'   # Bear: Jan 2018 - Des 2018
# Recovery: Jan 2019 - Okt 2019

periods = {
    'Bull (Nov-Dec 2017)' : df_returns.loc[df_returns.index <= BULL_END],
    'Bear (2018)'         : df_returns.loc[
                                (df_returns.index > BULL_END) &
                                (df_returns.index <= BEAR_END)],
    'Recovery (2019)'     : df_returns.loc[df_returns.index > BEAR_END],
    'Full Period'         : df_returns,
}

print('=== DATASET SUMMARY ===')
print(f'Assets  : {df_returns.columns.tolist()}')
print(f'Period  : {df_returns.index.min().date()} -> {df_returns.index.max().date()}')
print(f'Obs     : {len(df_returns)} trading days')
for pname, pdf in periods.items():
    print(f'{pname:<22}: {len(pdf):>4} days')
print(df_returns.describe().round(4))
\end{lstlisting}

%=============================================================================
\section{Sel 3 -- Helper Functions \& Risk Metrics}
%=============================================================================

\subsection{Deskripsi}
Mendefinisikan fungsi-fungsi utilitas untuk pemrosesan matriks korelasi menggunakan \textit{Random Matrix Theory} (RMT), pembangunan \textit{Minimum Spanning Tree} (MST), perhitungan sentralitas, serta berbagai metrik risiko dan uji statistik Diebold-Mariano.

\subsection{Kode}
\begin{lstlisting}[caption={Helper functions dan metrik risiko}]
def apply_rmt_filter(returns_data):
    T, N = returns_data.shape
    Q = T / N
    C = np.corrcoef(returns_data.T)
    eigenvalues, eigenvectors = eigh(C)
    # Balik urutan (terbesar ke terkecil)
    eigenvalues  = eigenvalues[::-1]
    eigenvectors = eigenvectors[:, ::-1]
    # Batas Marchenko-Pastur
    lambda_plus = 1 + (1/Q) + 2*np.sqrt(1/Q)
    # Nolkan nilai eigen yang dianggap noise
    Lambda_f = np.diag(np.where(eigenvalues > lambda_plus, eigenvalues, 0))
    return eigenvectors @ Lambda_f @ eigenvectors.T

def build_mst(correlation_matrix):
    # Pastikan korelasi berada dalam rentang [-1, 1] untuk menghindari NaN dalam sqrt
    correlation_matrix = np.clip(correlation_matrix, -1, 1)
    d = np.sqrt(2 - 2 * correlation_matrix)
    np.fill_diagonal(d, 0)
    return d

def compute_eigenvector_centrality(distance_matrix):
    adj = 1 / (distance_matrix + 1e-8)   # adjacency = 1/jarak
    np.fill_diagonal(adj, 0)
    _, vecs = eigh(adj)
    pev = np.abs(vecs[:, -1])             # eigenvector terbesar
    return pev / pev.sum()                # normalisasi

def get_assets_graph_diversify(returns_window, corr_threshold=0.4):
    corr_mat = returns_window.corr()
    G = nx.Graph()
    # Urutkan aset berdasarkan rata-rata return (terbesar duluan)
    assets = list(returns_window.mean().sort_values(ascending=False).index)
    G.add_nodes_from(assets)
    for i, a1 in enumerate(assets):
        for a2 in assets[i+1:]:
            if abs(corr_mat.loc[a1, a2]) > corr_threshold:
                G.add_edge(a1, a2)
    return list(nx.approximation.maximum_independent_set(G))

# Metrik Risiko
def calculate_var(returns, confidence=0.95):
    """Value at Risk pada level kepercayaan tertentu."""
    return np.percentile(returns, (1 - confidence) * 100)

def calculate_rachev_ratio(returns, alpha=0.10):
    """Rasio rata-rata tail atas terhadap tail bawah."""
    tu = np.percentile(returns, (1 - alpha) * 100)
    tl = np.percentile(returns, alpha * 100)
    cu = returns[returns >= tu].mean()
    cl = abs(returns[returns <= tl].mean())
    return cu / cl if cl > 0 else 0

def calculate_max_drawdown(cumulative_returns):
    """Drawdown maksimum dari nilai kumulatif portofolio."""
    running_max = np.maximum.accumulate(cumulative_returns)
    drawdown = (cumulative_returns - running_max) / (running_max + 1e-12)
    return drawdown.min()

def calculate_calmar_ratio(returns, cum_returns):
    """Calmar Ratio = Return Tahunan / |Max Drawdown|."""
    ann_ret = np.mean(returns) * 252
    mdd     = abs(calculate_max_drawdown(cum_returns))
    return ann_ret / mdd if mdd > 0 else 0

def diebold_mariano_test(returns_a, returns_b, h=1):
    T   = len(returns_a)
    d   = returns_a - returns_b
    d_mean = np.mean(d)
    def autocov(xi, xm, k):
        T_ = len(xi)
        return np.sum((xi[:T_-k] - xm) * (xi[k:] - xm)) / T_ if T_ > k else 0
    var_d = autocov(d, d_mean, 0)
    for i in range(1, h):
        var_d += 2 * autocov(d, d_mean, i)
    dm = d_mean / np.sqrt(abs(var_d) / T) if var_d != 0 else 0
    pv = 2 * (1 - stats.norm.cdf(abs(dm)))
    return dm, pv
\end{lstlisting}

%=============================================================================
\section{Sel 4 -- Graph Neural Network (GCN Architecture)}
%=============================================================================

\subsection{Deskripsi}
Arsitektur GCN dua lapis untuk menghasilkan \textit{node embeddings}. Korelasi 
antar aset diprediksi melalui \textit{cosine similarity} pada ruang laten. Bagian 
ini juga menyertakan fungsi untuk menyiapkan data graf dan melatih model.

\subsection{Kode}
\begin{lstlisting}[caption={Arsitektur CorrelationGNN dan pelatihan}]
class CorrelationGNN(nn.Module):
    def __init__(self, input_dim=4, hidden_dim=16, output_dim=8):
        super().__init__()
        self.conv1 = GCNConv(input_dim, hidden_dim)
        self.conv2 = GCNConv(hidden_dim, output_dim)

    def forward(self, x, edge_index):
        x = F.relu(self.conv1(x, edge_index))
        x = F.dropout(x, p=0.2, training=self.training)
        x = F.relu(self.conv2(x, edge_index))
        return x

    def predict_correlation(self, emb):
        n = emb.shape[0]
        C = torch.zeros(n, n)
        for i in range(n):
            for j in range(i+1, n):
                s = F.cosine_similarity(emb[i].unsqueeze(0), emb[j].unsqueeze(0))
                C[i, j] = C[j, i] = s
        for i in range(n): C[i,i] = 1.0
        return C

def create_graph_data(returns_window, corr_threshold=0.3):
    n = returns_window.shape[1]
    feats = np.stack([
        returns_window.mean().values,
        returns_window.std().values,
        returns_window.skew().values,
        returns_window.kurtosis().values
    ], axis=1)
    x = torch.FloatTensor(feats)
    corr = returns_window.corr().values
    edges = [[i, j] for i in range(n) for j in range(i+1, n) if abs(corr[i, j]) > corr_threshold]
    edges += [[j, i] for i, j in edges]
    if not edges:
        edges = [[i, j] for i in range(n) for j in range(n) if i != j]
    ei = torch.LongTensor(edges).t().contiguous()
    return Data(x=x, edge_index=ei, y=torch.FloatTensor(corr))

def train_gnn(model, data, epochs=30, lr=0.01):
    opt = torch.optim.Adam(model.parameters(), lr=lr)
    loss_fn = nn.MSELoss()
    model.train()
    for _ in range(epochs):
        opt.zero_grad()
        emb = model(data.x, data.edge_index)
        loss = loss_fn(model.predict_correlation(emb), data.y)
        loss.backward()
        opt.step()
    return model
\end{lstlisting}

%=============================================================================
\section{Sel 5 -- Portfolio Strategy Classes}
%=============================================================================

\subsection{Deskripsi}
Bagian ini mendefinisikan \textit{blueprint} utama untuk strategi portofolio melalui kelas basis \textit{PortfolioStrategy}. Berbagai strategi diimplementasikan sebagai turunan kelas ini, mulai dari strategi klasik hingga strategi berbasis graf dan GNN.

\subsection{Kode}
\begin{lstlisting}[caption={Implementasi berbagai strategi portofolio}]
class PortfolioStrategy:
    def __init__(self, name):
        self.name = name
    def get_weights(self, r):
        raise NotImplementedError

class EquallyWeighted(PortfolioStrategy):
    def get_weights(self, r):
        n = r.shape[1]
        return np.ones(n) / n

class ClassicalMarkowitz(PortfolioStrategy):
    def get_weights(self, r):
        n  = r.shape[1]
        mu = r.mean().values
        S  = r.cov().values
        f  = lambda w: w @ S @ w
        c  = [{'type': 'eq',  'fun': lambda w: np.sum(w) - 1},
              {'type': 'ineq','fun': lambda w: w @ mu - mu.mean()}]
        res = minimize(f, np.ones(n)/n, method='SLSQP',
                       bounds=[(0,1)]*n, constraints=c)
        return res.x if res.success else np.ones(n)/n

class GlassoMarkowitz(PortfolioStrategy):
    def __init__(self, name='Glasso Markowitz', alpha=0.01):
        super().__init__(name)
        self.alpha = alpha
    def get_weights(self, r):
        n  = r.shape[1]
        mu = r.mean().values
        try:
            from sklearn.covariance import GraphicalLassoCV
            g = GraphicalLassoCV(alphas=[self.alpha], cv=3)
            g.fit(r.values)
            S = g.covariance_
        except:
            S = r.cov().values
        f  = lambda w: w @ S @ w
        c  = [{'type': 'eq',  'fun': lambda w: np.sum(w) - 1},
              {'type': 'ineq','fun': lambda w: w @ mu - mu.mean()}]
        res = minimize(f, np.ones(n)/n, method='SLSQP',
                       bounds=[(0,1)]*n, constraints=c)
        return res.x if res.success else np.ones(n)/n

class NetworkMarkowitz(PortfolioStrategy):
    def __init__(self, name='Network Markowitz', gamma=0):
        super().__init__(name)
        self.gamma = gamma
    def get_weights(self, r):
        n    = r.shape[1]
        mu   = r.mean().values
        sig  = r.std().values
        Cf   = apply_rmt_filter(r)
        dm   = build_mst(Cf)
        cent = compute_eigenvector_centrality(dm)
        Sf   = np.outer(sig, sig) * Cf
        f    = lambda w: w @ Sf @ w + self.gamma * np.sum(cent * w)
        c    = [{'type': 'eq',  'fun': lambda w: np.sum(w) - 1},
                {'type': 'ineq','fun': lambda w: w @ mu - mu.mean()}]
        res  = minimize(f, np.ones(n)/n, method='SLSQP',
                        bounds=[(0,1)]*n, constraints=c)
        return res.x if res.success else np.ones(n)/n

class GraphDiversification(PortfolioStrategy):
    def __init__(self, name='Graph Diversification', corr_threshold=0.4):
        super().__init__(name)
        self.corr_threshold = corr_threshold
    def get_weights(self, r):
        n      = r.shape[1]
        assets = r.columns.tolist()
        sel    = get_assets_graph_diversify(r, self.corr_threshold)
        w      = np.zeros(n)
        if sel:
            for a in sel: w[assets.index(a)] = 1.0 / len(sel)
        else:
            w = np.ones(n) / n
        return w

class GNNGraphDiversification(PortfolioStrategy):
    def __init__(self, name='GNN GD', train_epochs=30):
        super().__init__(name)
        self.train_epochs = train_epochs
        self.gnn_model    = CorrelationGNN()
    def get_weights(self, r):
        n = r.shape[1]; assets = r.columns.tolist()
        data = create_graph_data(r)
        train_gnn(self.gnn_model, data, epochs=self.train_epochs)
        self.gnn_model.eval()
        with torch.no_grad():
            emb = self.gnn_model(data.x, data.edge_index)
            pc  = self.gnn_model.predict_correlation(emb).numpy()
        up = pc[np.triu_indices(n, k=1)]
        thr = np.clip(np.mean(np.abs(up)) + 0.5*np.std(np.abs(up)), 0.3, 0.7)
        G = nx.Graph(); asorted = list(r.mean().sort_values(ascending=False).index)
        G.add_nodes_from(asorted)
        for i, a1 in enumerate(asorted):
            for j, a2 in enumerate(asorted):
                if i < j:
                    i1, i2 = assets.index(a1), assets.index(a2)
                    if abs(pc[i1, i2]) > thr: G.add_edge(a1, a2)
        sel = list(nx.approximation.maximum_independent_set(G))
        w = np.zeros(n)
        if sel:
            for a in sel: w[assets.index(a)] = 1.0 / len(sel)
        else:
            w = np.ones(n) / n
        return w
\end{lstlisting}

%=============================================================================
\section{Sel 6 -- Backtesting \& Performance Metrics}
%=============================================================================

\subsection{Deskripsi}
Sistem backtesting melakukan iterasi \textit{rolling window} dengan memperhitungkan biaya transaksi. Hasil simulasi kemudian diproses melalui fungsi metrik untuk mendapatkan data statistik performa portofolio.

\subsection{Kode}
\begin{lstlisting}[caption={Framework backtesting dan metrik}]
def backtest_strategy(strategy, df_ret, window_size=120, rebalance_freq=7, tc=0.001):
    port_returns = []; dates = []
    for i in range(window_size, len(df_ret), rebalance_freq):
        train = df_ret.iloc[i-window_size:i]
        w = strategy.get_weights(train)
        test_end = min(i + rebalance_freq, len(df_ret))
        test = df_ret.iloc[i:test_end]
        for j in range(len(test)):
            dr = np.dot(w, test.iloc[j].values)
            if j == 0 and port_returns: dr -= tc
            port_returns.append(dr)
            dates.append(test.index[j])
    res = pd.DataFrame({'date': dates, 'return': port_returns})
    res['cumulative_return'] = (1 + res['return']).cumprod()
    return {
        'strategy': strategy.name,
        'returns': np.array(port_returns),
        'cumulative_returns': res['cumulative_return'].values,
        'dates': dates,
        'results_df': res
    }

def calculate_metrics(result):
    r = result['returns']; cr = result['cumulative_returns']
    ann_r = np.mean(r) * 252; ann_v = np.std(r) * np.sqrt(252)
    sharpe = ann_r / ann_v if ann_v > 0 else 0
    return {
        'Strategy': result['strategy'],
        'Total Return (%)': round((cr[-1] - 1) * 100, 2),
        'Annual Return (%)': round(ann_r * 100, 2),
        'Annual Vol (%)': round(ann_v * 100, 2),
        'Sharpe Ratio': round(sharpe, 4),
        'VaR 95% (%)': round(calculate_var(r, 0.95) * 100, 4),
        'Rachev Ratio': round(calculate_rachev_ratio(r, 0.10), 4),
        'Max Drawdown (%)': round(calculate_max_drawdown(cr) * 100, 2),
        'Calmar Ratio': round(calculate_calmar_ratio(r, cr), 4),
    }

# Evaluasi Metrik dan Highlighting
metrics_df = pd.DataFrame(
    [calculate_metrics(r) for r in results.values()]
).set_index('Strategy')

def highlight_best(col):
    """Highlight nilai terbaik per kolom (hijau untuk performa terbaik)."""
    # Kolom di mana nilai terkecil adalah yang terbaik (risiko)
    is_min = col.name in ['Annual Vol (%)', 'VaR 95% (%)', 'Max Drawdown (%)']
    best   = col.min() if is_min else col.max()
    return ['background-color: #d4edda; font-weight: bold'
            if v == best else '' for v in col]

print('PERFORMANCE METRICS (Full Period)')
print('='*100)
print(metrics_df.to_string())
print('='*100)

# Menampilkan dataframe dengan styling (dalam lingkungan Notebook)
metrics_df.style.apply(highlight_best)
\end{lstlisting}

\subsection{Visualisasi Hasil}
Selain tabel metrik, visualisasi kumulatif return digunakan untuk membandingkan stabilitas portofolio dari waktu ke waktu.

\begin{lstlisting}[caption={Plotting Cumulative Returns}]
plt.figure(figsize=(12, 6))
for sname, res in results.items():
    plt.plot(res['dates'], res['cumulative_returns'], label=sname)
plt.title('Cumulative Returns Comparison')
plt.xlabel('Date'); plt.ylabel('Portfolio Value')
plt.legend(); plt.show()
\end{lstlisting}

\subsection{Validasi Statistik (Diebold-Mariano Test)}
Untuk memastikan perbedaan performa antar strategi signifikan secara statistik, dilakukan uji Diebold-Mariano secara berpasangan (\textit{pairwise}).

\begin{lstlisting}[caption={Komputasi Matriks DM Test}]
strategy_names = list(results.keys())
dm_stat_mat = pd.DataFrame(index=strategy_names, columns=strategy_names, dtype=float)
dm_pval_mat = pd.DataFrame(index=strategy_names, columns=strategy_names, dtype=float)

for i, s1 in enumerate(strategy_names):
    for j, s2 in enumerate(strategy_names):
        if i == j:
            dm_stat_mat.loc[s1, s2] = 0.0
            dm_pval_mat.loc[s1, s2] = 1.0
        else:
            dm, pv = diebold_mariano_test(results[s1]['returns'], results[s2]['returns'])
            dm_stat_mat.loc[s1, s2] = round(dm, 3)
            dm_pval_mat.loc[s1, s2] = round(pv, 4)

sig_matrix = dm_pval_mat < 0.05
print('DM Test P-Value Matrix (p < 0.05 = significant)')
print(dm_pval_mat.to_string())
print('\nDM Test Statistic Matrix (positive = row outperforms column)')
print(dm_stat_mat.to_string())

# Visualisasi Heatmap DM Test
short_names = {
    'Equally Weighted': 'EW',
    'Classical Markowitz': 'CMV',
    'Glasso Markowitz': 'Glasso',
    'Network Markowitz (gamma=0)': 'Net-MV(0)',
    'Network Markowitz (gamma=1.0)': 'Net-MV(1)',
    'Graph Divers. (theta=0.4)': 'GD(0.4)',
    'Graph Divers. (theta=0.5)': 'GD(0.5)',
    'GNN Graph Divers. (30 ep)': 'GNN-30',
    'GNN Graph Divers. (50 ep)': 'GNN-50',
}
dm_stat_r = dm_stat_mat.rename(index=short_names, columns=short_names).astype(float)
pval_r    = dm_pval_mat.rename(index=short_names, columns=short_names).astype(float)

fig, ax = plt.subplots(figsize=(11, 8))
mask_diag = np.eye(len(dm_stat_r), dtype=bool)
sns.heatmap(dm_stat_r, annot=True, fmt='.2f', center=0, cmap='RdYlGn',
            linewidths=0.5, ax=ax, mask=mask_diag,
            cbar_kws={'label': 'DM Statistic (positive = row outperforms col)'})

for i in range(len(dm_stat_r)):
    for j in range(len(dm_stat_r.columns)):
        if i != j and pval_r.iloc[i, j] < 0.05:
            ax.text(j+0.85, i+0.15, '*', color='black', fontsize=14, fontweight='bold')

plt.title('Diebold-Mariano Test Statistic Heatmap\n(* = significant at p<0.05)', fontsize=12)
plt.tight_layout()
plt.savefig('dm_test_heatmap.png', dpi=150); plt.show()
\end{lstlisting}



%=============================================================================
\section{Kesimpulan Arsitektur GNN}
%=============================================================================

\begin{itemize}
    \item \textbf{Input Dimensionality}: 4 fitur node (Mean, Std, Skew, Kurt).
    \item \textbf{Output Dimensionality}: 8-dim latent space (Embedding).
    \item \textbf{Adaptabilitas}: Threshold graf dihitung ulang setiap rebalancing minggu melalui model yang dilatih online.
    \item \textbf{Optimasi}: Adam Optimizer dengan MSE Loss untuk meminimalkan selisih korelasi laten dan empiris.
\end{itemize}

\end{document}


